% =====================================================================
\chapter{Soluciones de los ejercicios}
\label{ap:soluciones}
% =====================================================================

Cada solución se presenta en cinco partes: datos, planteamiento, desarrollo, resultado y comprobación. En los pasos intermedios se redondea a cuatro cifras significativas salvo que se indique otra cosa; los resultados exactos se dan como fracción cuando existe.

% ---------------------------------------------------------------------
\section{Tema 1}

\begin{solucion}{eje:1-cuartica}
\datos $f(x)=x^{4}$, $f'(x)=4x^{3}$, $f''(x)=12x^{2}$.
\planteamiento usar las caracterizaciones de segundo orden de la convexidad y de la convexidad fuerte (\cref{def:fuerte}) y la definición de $L$-suavidad (\cref{def:suave}).
\pasos
\begin{enumerate}[label=\arabic*.]
  \item (a) $f''(x)=12x^{2}\ge 0$ para todo $x$: $f$ es convexa.
  \item (b) La convexidad fuerte exigiría $f''(x)\ge\mu>0$ para todo $x$, pero $f''(0)=0$. No es fuertemente convexa.
  \item (c) $\abs{f'(x)-f'(0)}=4\abs{x}^{3}$, y el cociente $4\abs{x}^{3}/\abs{x}=4x^{2}$ no está acotado en $\R$. No existe $L$ global.
  \item Unicidad: $f(x)>0=f(0)$ si $x\ne 0$, luego $0$ es el único minimizador. La convexidad fuerte es suficiente para la unicidad, no necesaria: basta la convexidad estricta, que $f$ sí tiene.
\end{enumerate}
\resultado \textbf{convexa (incluso estrictamente), no fuertemente convexa y no $L$-suave en $\R$}; no hay contradicción.
\comprobacion cerca del mínimo la curvatura es nula, así que el descenso de gradiente con paso fijo converge solo sublinealmente: con $\eta$ pequeño, $x_{k+1}=x_k(1-4\eta x_k^{2})$, cuyo factor tiende a~1.
\end{solucion}

\begin{solucion}{eje:1-silla}
\datos $f(x,y)=x^{2}-y^{2}$; punto inicial $(1,0)$ o $(1,\delta)$; $\eta=0{,}1$; $\delta=10^{-6}$.
\planteamiento clasificar con el hessiano (\cref{prop:clasificacion}) y seguir la iteración por coordenadas (\cref{ej:silla}).
\pasos
\begin{enumerate}[label=\arabic*.]
  \item $\nabla f=(2x,-2y)$ se anula en $(0,0)$.
  \item $\nabla^{2}f=\diag(2,-2)$, con autovalores de signos opuestos: \emph{punto de silla}.
  \item Iteración: $x_{k+1}=(1-2\eta)x_k$, $y_{k+1}=(1+2\eta)y_k$.
  \item Con $y_0=0$: $y_k=0$ para todo $k$. Con $0<\eta<1$, $\abs{1-2\eta}<1$ y $x_k\to 0$: el método converge a la silla.
  \item Con $y_0=\delta$: $y_k=(1+2\eta)^{k}\delta=1{,}2^{k}\cdot 10^{-6}$. Para $k=50$, $1{,}2^{50}\approx 9100$ e $y_{50}\approx 9{,}1\cdot 10^{-3}$.
\end{enumerate}
\resultado \textbf{$(0,0)$ es un punto de silla}; desde $(1,0)$ el método converge a él porque la coordenada inestable es exactamente nula; con $\delta=10^{-6}$, $y_{50}\approx 9{,}1\cdot 10^{-3}$ y sigue creciendo.
\comprobacion $\ln(1{,}2^{50})=50\cdot 0{,}1823=9{,}116$ y $e^{9{,}116}\approx 9100$.
\end{solucion}

\begin{solucion}{eje:1-cuadratica}
\datos $Q=\begin{pmatrix}2&1\\1&2\end{pmatrix}$, $c=(-3,0)\T$, $b=1$.
\planteamiento $Q$ tiene autovalores $1$ y $3$, luego $Q\succ 0$ y se aplica la \cref{prop:min-cuadratica}.
\pasos
\begin{enumerate}[label=\arabic*.]
  \item $\det Q=4-1=3$; $Q^{-1}=\frac13\begin{pmatrix}2&-1\\-1&2\end{pmatrix}$.
  \item $x^{\ast}=-Q^{-1}c=-\frac13\bigl(2\cdot(-3)+(-1)\cdot 0;\ (-1)\cdot(-3)+2\cdot 0\bigr)=-\frac13(-6;\,3)=(2;\,-1)$.
  \item $c\T Q^{-1}c=-c\T x^{\ast}=-\bigl((-3)\cdot 2+0\cdot(-1)\bigr)=6$.
  \item $h^{\ast}=b-\frac12c\T Q^{-1}c=1-3=-2$.
\end{enumerate}
\resultado \textbf{$x^{\ast}=(2,-1)$ y $h(x^{\ast})=-2$.}
\comprobacion $\frac12x^{\ast\top}Qx^{\ast}=\frac12(2\cdot 4+2\cdot 2\cdot(-1)+2\cdot 1)=\frac12\cdot 6=3$; $x^{\ast\top}c=-6$; $h=3-6+1=-2$. Además $Qx^{\ast}+c=(4-1-3;\ 2-2+0)=(0;\,0)$.
\end{solucion}

\begin{solucion}{eje:1-constantes}
\datos $f(x,y)=x^{2}+2xy+5y^{2}$.
\planteamiento $f=\frac12\theta\T Q\theta$ con $Q$ el hessiano; $\mu=\lambda_{\min}(Q)$ y $L=\lambda_{\max}(Q)$ (\cref{ej:cuadratica}).
\pasos
\begin{enumerate}[label=\arabic*.]
  \item $\partial^{2}f/\partial x^{2}=2$, $\partial^{2}f/\partial x\partial y=2$, $\partial^{2}f/\partial y^{2}=10$: $Q=\begin{pmatrix}2&2\\2&10\end{pmatrix}$.
  \item Traza $12$ y determinante $20-4=16$. Ecuación característica: $\lambda^{2}-12\lambda+16=0$.
  \item $\lambda=6\pm\sqrt{36-16}=6\pm 2\sqrt5$: $\lambda_{\min}=6-2\sqrt5\approx 1{,}528$ y $\lambda_{\max}=6+2\sqrt5\approx 10{,}47$.
  \item Ambos son positivos, luego $Q\succ 0$ y $f$ es $\mu$-fuertemente convexa con $\mu\approx 1{,}528$, y $L$-suave con $L\approx 10{,}47$.
  \item $\kappa=\dfrac{6+2\sqrt5}{6-2\sqrt5}=\dfrac{(3+\sqrt5)^{2}}{9-5}=\dfrac{14+6\sqrt5}{4}=3{,}5+1{,}5\sqrt5\approx 6{,}854$.
\end{enumerate}
\resultado \textbf{$\mu=6-2\sqrt5\approx 1{,}528$, $L=6+2\sqrt5\approx 10{,}47$ y $\kappa\approx 6{,}854$.}
\comprobacion $\mu+L=12$ (traza) y $\mu L=36-20=16$ (determinante).
\end{solucion}

\begin{solucion}{eje:1-maximo}
\datos $f,g$ convexas en $\R^{n}$; $h=\max\{f,g\}$.
\planteamiento aplicar la definición de convexidad a $h$ y acotar cada función por $h$.
\pasos
\begin{enumerate}[label=\arabic*.]
  \item Sean $x,y$ y $t\in[0,1]$, y $z=tx+(1-t)y$. Por la convexidad de $f$, $f(z)\le tf(x)+(1-t)f(y)$.
  \item Como $f\le h$ y $t,1-t\ge 0$, $tf(x)+(1-t)f(y)\le th(x)+(1-t)h(y)$. Luego $f(z)\le th(x)+(1-t)h(y)$.
  \item Del mismo modo, $g(z)\le th(x)+(1-t)h(y)$.
  \item Por tanto, $h(z)=\max\{f(z),g(z)\}\le th(x)+(1-t)h(y)$: $h$ es convexa.
  \item (b) Contraejemplo en $\R$: $f(x)=(x-1)^{2}$ y $g(x)=(x+1)^{2}$. Sea $m=\min\{f,g\}$. Entonces $m(1)=m(-1)=0$ y $m(0)=1$.
  \item Con $x=1$, $y=-1$ y $t=\frac12$: $m(0)=1>\frac12m(1)+\frac12m(-1)=0$. La desigualdad de convexidad falla.
\end{enumerate}
\resultado \textbf{el máximo de funciones convexas es convexo; el mínimo, en general, no.}
\comprobacion la gráfica de $m$ tiene forma de «W», con dos mínimos locales en $\pm 1$; una función convexa no puede tener dos mínimos locales separados por un punto de valor mayor (\cref{teo:opt-convexa}).
\end{solucion}

\begin{solucion}{eje:1-fuerte}
\datos $f$ diferenciable, $\mu$-fuertemente convexa, con minimizador $\theta^{\ast}$.
\planteamiento aplicar la \cref{def:fuerte} en el punto $\theta^{\ast}$, donde $\nabla f(\theta^{\ast})=0$ (\cref{teo:fermat}).
\pasos
\begin{enumerate}[label=\arabic*.]
  \item (a) Con $x=\theta^{\ast}$ e $y=\theta$: $f(\theta)\ge f(\theta^{\ast})+0+\frac{\mu}{2}\norm{\theta-\theta^{\ast}}^{2}$.
  \item (b) Si $\theta'$ es otro minimizador, $f(\theta')=f^{\ast}$ y (a) da $0\ge\frac{\mu}{2}\norm{\theta'-\theta^{\ast}}^{2}$, luego $\theta'=\theta^{\ast}$.
  \item (c) Se escribe la \cref{def:fuerte} con $x=\theta$, $y=\theta^{\ast}$: $f^{\ast}\ge f(\theta)+\nabla f(\theta)\T(\theta^{\ast}-\theta)+\frac{\mu}{2}\norm{\theta-\theta^{\ast}}^{2}$.
  \item Se suma con la desigualdad de (a), $f(\theta)\ge f^{\ast}+\frac{\mu}{2}\norm{\theta-\theta^{\ast}}^{2}$: queda $0\ge\nabla f(\theta)\T(\theta^{\ast}-\theta)+\mu\norm{\theta-\theta^{\ast}}^{2}$, es decir, $\mu\norm{\theta-\theta^{\ast}}^{2}\le\nabla f(\theta)\T(\theta-\theta^{\ast})$.
  \item Por Cauchy--Schwarz, $\nabla f(\theta)\T(\theta-\theta^{\ast})\le\norm{\nabla f(\theta)}\,\norm{\theta-\theta^{\ast}}$. Dividiendo por $\norm{\theta-\theta^{\ast}}$ (si es no nulo; si es nulo, la desigualdad es trivial), $\mu\norm{\theta-\theta^{\ast}}\le\norm{\nabla f(\theta)}$.
\end{enumerate}
\resultado \textbf{se cumplen (a), (b) y (c)}: la función crece al menos cuadráticamente lejos del óptimo, y el gradiente, al menos linealmente.
\comprobacion para $f(\theta)=\frac{\mu}{2}\norm{\theta}^{2}$, con $\theta^{\ast}=0$, las tres desigualdades son igualdades: $f-f^{\ast}=\frac{\mu}{2}\norm{\theta}^{2}$ y $\norm{\nabla f}=\mu\norm{\theta}$.
\end{solucion}

% ---------------------------------------------------------------------
\section{Tema 2}

\begin{solucion}{eje:2-escala}
\datos $f$ con minimizador $\theta^{\ast}$, constantes $\mu$, $L$; se minimiza $cf$ con $c>0$.
\planteamiento $\nabla(cf)=c\nabla f$ y $\nabla^{2}(cf)=c\nabla^{2}f$.
\pasos
\begin{enumerate}[label=\arabic*.]
  \item Minimizador: $cf(\theta)\ge cf(\theta^{\ast})$ para todo $\theta$, luego el minimizador no cambia.
  \item Las constantes pasan a ser $c\mu$ y $cL$: el paso máximo estable, $2/L$ en una cuadrática, pasa a ser $2/(cL)$.
  \item $\kappa=cL/(c\mu)=L/\mu$: el número de condición no cambia.
  \item Con paso $\eta'$ sobre $cf$: $\theta_{k+1}=\theta_k-\eta'c\nabla f(\theta_k)$. Coincide con el descenso sobre $f$ con paso $\eta$ si y solo si $\eta'=\eta/c$.
\end{enumerate}
\resultado \textbf{el minimizador y $\kappa$ no cambian; el paso máximo se divide por $c$; con $\eta/c$ se reproducen los mismos iterados.}
\comprobacion contrasta con Adam, cuyos iterados no cambian aunque se mantenga el mismo $\eta$ (\cref{prop:escala}).
\end{solucion}

\begin{solucion}{eje:2-1d}
\datos $f(x)=3x^{2}$, $x_0=2$, $\eta=0{,}1$.
\planteamiento $f=\frac{a}{2}x^{2}$ con $a=6$; cada paso multiplica $x$ por $1-\eta a$ (\cref{ej:gd-1d}).
\pasos
\begin{enumerate}[label=\arabic*.]
  \item $f'(x)=6x$; factor $1-0{,}1\cdot 6=0{,}4$.
  \item $x_1=0{,}4\cdot 2=0{,}8$; $x_2=0{,}4\cdot 0{,}8=0{,}32$; $x_3=0{,}4\cdot 0{,}32=0{,}128$.
  \item Convergencia: $\abs{1-6\eta}<1\iff 0<\eta<2/6=1/3$.
\end{enumerate}
\resultado \textbf{$x_1=0{,}8$, $x_2=0{,}32$, $x_3=0{,}128$; converge si y solo si $0<\eta<1/3\approx 0{,}3333$.} Con $\eta=1/3$ el factor es $-1$ y el iterado oscila entre $\pm 2$.
\comprobacion $x_3=2\cdot 0{,}4^{3}=2\cdot 0{,}064=0{,}128$. Con $\eta=0{,}4$ el factor sería $1-2{,}4=-1{,}4$, de módulo mayor que~1.
\end{solucion}

\begin{solucion}{eje:2-2d}
\datos $f(x,y)=x^{2}+5y^{2}$, $\theta_0=(1,1)$, $\eta=1/6$.
\planteamiento $Q=\diag(2,10)$; cada coordenada se multiplica por $1-\eta\lambda_i$ (\cref{teo:gd-cuadratica}).
\pasos
\begin{enumerate}[label=\arabic*.]
  \item Factor en $x$: $1-\frac16\cdot 2=\frac23$. Factor en $y$: $1-\frac16\cdot 10=-\frac23$.
  \item $\theta_1=\bigl(\frac23;\,-\frac23\bigr)$, $\theta_2=\bigl(\frac49;\,\frac49\bigr)$, $\theta_3=\bigl(\frac{8}{27};\,-\frac{8}{27}\bigr)\approx(0{,}2963;\,-0{,}2963)$.
  \item Normas: $\norm{\theta_0}=\sqrt2\approx 1{,}414$; $\norm{\theta_1}=\frac23\sqrt2\approx 0{,}9428$; $\norm{\theta_2}=\frac49\sqrt2\approx 0{,}6285$; $\norm{\theta_3}=\frac{8}{27}\sqrt2\approx 0{,}4190$.
  \item Cociente entre normas consecutivas: $\frac23$.
  \item Teorema: $\mu=2$, $L=10$, $\kappa=5$; $\eta^{\ast}=2/12=1/6$ y $\rho=(5-1)/(5+1)=\frac23$.
\end{enumerate}
\resultado \textbf{$\theta_3=(8/27;\,-8/27)$ y la norma se multiplica por $2/3$ en cada paso}, que es $(\kappa-1)/(\kappa+1)$ porque $\eta=1/6$ es el paso óptimo.
\comprobacion $f(\theta_3)=6\cdot(8/27)^{2}=384/729\approx 0{,}5267$ y $f(\theta_3)/f(\theta_0)=0{,}5267/6\approx 0{,}08779=(2/3)^{6}$. La coordenada $y$ cambia de signo en cada paso, como en la \cref{fig:zigzag}.
\end{solucion}

\begin{solucion}{eje:2-valle}
\datos $f=\frac12(x^{2}+25y^{2})$, $Q=\diag(1,25)$, $\theta_0=(1,1)$, $\epsilon=10^{-3}$.
\planteamiento \cref{teo:gd-cuadratica} y \cref{cor:iter-cuadratica}: $\eta^{\ast}=2/(\mu+L)$, $\rho=(\kappa-1)/(\kappa+1)$, $k=\lceil\ln\epsilon/\ln\rho\rceil$.
\pasos
\begin{enumerate}[label=\arabic*.]
  \item $\mu=1$, $L=25$, $\kappa=25$.
  \item $\eta^{\ast}=2/26=1/13\approx 0{,}07692$; $\rho=24/26=12/13\approx 0{,}9231$.
  \item $\ln 10^{-3}=-6{,}908$; $\ln(12/13)=-0{,}08004$; $6{,}908/0{,}08004\approx 86{,}3$, luego \textbf{87} iteraciones.
  \item Momentum: $\sqrt\kappa=5$, $\rho=4/6=2/3$, $\ln(2/3)=-0{,}4055$; $6{,}908/0{,}4055\approx 17{,}04$, luego \textbf{18} iteraciones.
  \item Factores con $\eta^{\ast}$: $1-\frac{1}{13}=\frac{12}{13}$ en $x$ y $1-\frac{25}{13}=-\frac{12}{13}$ en $y$.
  \item $\theta_1=\bigl(\frac{12}{13};\,-\frac{12}{13}\bigr)\approx(0{,}9231;\,-0{,}9231)$ y $\theta_2=\bigl(\frac{144}{169};\,\frac{144}{169}\bigr)\approx(0{,}8521;\,0{,}8521)$.
\end{enumerate}
\resultado \textbf{$\kappa=25$, $\eta^{\ast}=1/13$, $\rho=12/13$; 87 iteraciones con descenso de gradiente frente a 18 con momentum.}
\comprobacion $\norm{\theta_1}/\norm{\theta_0}=12/13=\rho$. El cociente de iteraciones, $87/18\approx 4{,}8$, es próximo a $\sqrt\kappa=5$.
\end{solucion}

\begin{solucion}{eje:2-paso-largo}
\datos $F(x,y)=y^{2}+(y-x)^{2}$, $(x_0,y_0)=(1,1)$, pasos $s=1/2$ y $s=1/4$.
\planteamiento calcular el gradiente, iterar y analizar la estabilidad con $s<2/\lambda_{\max}$ (\cref{teo:gd-cuadratica}).
\pasos
\begin{enumerate}[label=\arabic*.]
  \item $\partial F/\partial x=-2(y-x)$; $\partial F/\partial y=2y+2(y-x)=4y-2x$. En $(1,1)$: $\nabla F=(0;\,2)$.
  \item $s=\frac12$. Paso~1: $(1,1)-\frac12(0,2)=(1,0)$, $F=0+1=1$.
  \item Paso~2: $\nabla F(1,0)=(2;\,-2)$; $(1,0)-\frac12(2,-2)=(0,1)$, $F=1+1=2$.
  \item Paso~3: $\nabla F(0,1)=(-2;\,4)$; $(0,1)-\frac12(-2,4)=(1,-1)$, $F=1+4=5$. $F$ crece: el método diverge.
  \item (b) $\nabla^{2}F=\begin{pmatrix}2&-2\\-2&4\end{pmatrix}$, traza 6, determinante 4: $\lambda=3\pm\sqrt5\approx 0{,}7639$ y $5{,}236$. Límite: $2/5{,}236\approx 0{,}3820<0{,}5$. El factor de la dirección rígida es $1-0{,}5\cdot 5{,}236=-1{,}618$, de módulo mayor que~1.
  \item (c) $s=\frac14$. Paso~1: $(1,1)-\frac14(0,2)=(1;\,0{,}5)$, $F=0{,}25+0{,}25=0{,}5$.
  \item Paso~2: $\nabla F(1;\,0{,}5)=(1;\,0)$; nuevo punto $(0{,}75;\,0{,}5)$, $F=0{,}25+0{,}0625=0{,}3125$.
  \item Paso~3: $\nabla F(0{,}75;\,0{,}5)=(0{,}5;\,0{,}5)$; nuevo punto $(0{,}625;\,0{,}375)$, $F=0{,}1406+0{,}0625=0{,}2031$.
\end{enumerate}
\resultado \textbf{con $s=1/2$, $F=1$; $2$; $5$: diverge. Con $s=1/4$, $F=0{,}5$; $0{,}3125$; $0{,}2031$: converge.}
\comprobacion con $s=\frac14$, los factores son $1-0{,}25\cdot 0{,}7639=0{,}8090$ y $1-0{,}25\cdot 5{,}236=-0{,}3090$, ambos de módulo menor que~1.
\end{solucion}

\begin{solucion}{eje:2-rosenbrock}
\datos $\nabla^{2}f(1,1)=\begin{pmatrix}802&-400\\-400&200\end{pmatrix}$, $\eta=10^{-3}$, $\epsilon=10^{-3}$.
\planteamiento autovalores de una matriz $2\times 2$ con la traza y el determinante; $\kappa=\lambda_{\max}/\lambda_{\min}$, paso máximo $2/\lambda_{\max}$.
\pasos
\begin{enumerate}[label=\arabic*.]
  \item Traza: $802+200=1002$. Determinante: $802\cdot 200-400^{2}=160\,400-160\,000=400$.
  \item $\lambda^{2}-1002\lambda+400=0$; discriminante $1002^{2}-1600=1\,002\,404$, de raíz $\approx 1001{,}2$.
  \item $\lambda_{\max}=(1002+1001{,}2)/2\approx 1001{,}6$. Para el menor se usa el producto, que evita restar números casi iguales: $\lambda_{\min}=400/1001{,}6\approx 0{,}3994$.
  \item $\kappa\approx 1001{,}6/0{,}3994\approx 2508$.
  \item Paso máximo: $2/1001{,}6\approx 1{,}997\cdot 10^{-3}$.
  \item Con $\eta=10^{-3}$, el factor en la dirección plana es $1-10^{-3}\cdot 0{,}3994=0{,}9996$. Iteraciones: $\ln 10^{-3}/\ln 0{,}9996\approx -6{,}908/(-3{,}995\cdot 10^{-4})\approx 1{,}73\cdot 10^{4}$.
\end{enumerate}
\resultado \textbf{$\lambda\approx 0{,}3994$ y $1001{,}6$; $\kappa\approx 2508$; paso máximo $\approx 0{,}002$; unas $17\,300$ iteraciones} para reducir el error en $10^{-3}$ en la dirección plana.
\comprobacion $\lambda_{\min}+\lambda_{\max}\approx 1002$ (traza). Es coherente con que el descenso de gradiente no llegue a $f<10^{-4}$ en 5000 iteraciones (\cref{ej:rosenbrock}).
\end{solucion}

\begin{solucion}{eje:2-cuadratica}
\datos $f(\theta)=\frac12\theta\T Q\theta$, $Q$ simétrica con autovalores $0<\mu=\lambda_1\le\dots\le\lambda_n=L$.
\planteamiento diagonalizar la iteración lineal $\theta_{k+1}=(I-\eta Q)\theta_k$ y estudiar $\rho(\eta)=\max_i\abs{1-\eta\lambda_i}$.
\pasos
\begin{enumerate}[label=\arabic*.]
  \item $\nabla f(\theta)=Q\theta$ por ser $Q$ simétrica (módulo~4), luego $\theta_{k+1}=(I-\eta Q)\theta_k$.
  \item Teorema espectral: $Q=U\Lambda U\T$ con $U$ ortogonal; $I-\eta Q=U(I-\eta\Lambda)U\T$, de autovalores $1-\eta\lambda_i$.
  \item Con $z=U\T\theta$, $z^{(i)}_k=(1-\eta\lambda_i)^{k}z^{(i)}_0$ y $\norm{\theta_k}=\norm{z_k}\le\rho(\eta)^{k}\norm{\theta_0}$.
  \item Si $0<\eta<2/L$: para todo $i$, $0<\eta\lambda_i\le\eta L<2$, luego $-1<1-\eta\lambda_i<1$ y $\rho(\eta)<1$: $\theta_k\to 0$.
  \item Extremos: $\varphi(\lambda)=\abs{1-\eta\lambda}$ es el valor absoluto de una función afín, luego convexa. Todo $\lambda\in[\mu,L]$ es $\lambda=t\mu+(1-t)L$ y $\varphi(\lambda)\le t\varphi(\mu)+(1-t)\varphi(L)\le\max\{\varphi(\mu),\varphi(L)\}$. Así, $\rho(\eta)=\max\{\abs{1-\eta\mu},\abs{1-\eta L}\}$.
  \item Óptimo: para $\eta\le 2/(\mu+L)$ el máximo es $1-\eta\mu$, decreciente en $\eta$; para $\eta\ge 2/(\mu+L)$ es $\eta L-1$, creciente. En efecto, $1-\eta\mu\ge\eta L-1\iff\eta\le 2/(\mu+L)$. El mínimo de $\rho$ está, por tanto, en $\eta^{\ast}=2/(\mu+L)$, donde ambos valores se igualan.
  \item $\rho(\eta^{\ast})=1-\frac{2\mu}{\mu+L}=\frac{L-\mu}{L+\mu}=\frac{\kappa-1}{\kappa+1}$.
\end{enumerate}
\resultado \textbf{el método converge para todo $\eta\in(0,2/L)$, y el paso óptimo $2/(\mu+L)$ da el factor $(\kappa-1)/(\kappa+1)$}. Igualar los extremos es óptimo porque cualquier otro $\eta$ empeora uno de los dos.
\comprobacion con $Q=\diag(2,10)$: $\eta^{\ast}=1/6$ y $\rho=2/3$, el valor observado en el \cref{eje:2-2d}.
\end{solucion}

% ---------------------------------------------------------------------
\section{Tema 3}

\begin{solucion}{eje:3-motivos}
\datos un corpus de billones de tokens y la regla $\theta\leftarrow\theta-\eta\nabla J(\theta)$, con $J$ la media sobre todo el corpus.
\planteamiento comparar el coste de un paso exacto con el de uno de minibatch y considerar el efecto del ruido.
\pasos
\begin{enumerate}[label=\arabic*.]
  \item \emph{Coste.} El gradiente exacto exige recorrer todo el corpus, con propagación y retropropagación, para dar un solo paso: un único paso costaría tanto como una época completa.
  \item \emph{Rendimiento por cálculo.} Por el \cref{teo:minibatch}, un minibatch de $B$ ejemplos da un estimador insesgado de varianza $\sigma^{2}/B$. Con el mismo cálculo que un paso exacto se dan $N/B$ pasos aproximados, que avanzan mucho más.
  \item \emph{Ruido útil.} El ruido ayuda a abandonar puntos de silla y se asocia a mínimos más anchos (\cref{obs:paisaje}).
\end{enumerate}
\resultado \textbf{se usa SGD con minibatch porque el gradiente exacto es inasumible y porque los pasos ruidosos son más eficientes.}
\comprobacion GPT-3 usó lotes de $3{,}2\cdot 10^{6}$ tokens~\cite[tabla~2.1]{gpt3}, una fracción ínfima de los cientos de miles de millones de tokens con que se entrenó.
\end{solucion}

\begin{solucion}{eje:3-minibatch}
\datos gradientes individuales $(-2,0,1,5)$, $N=4$, $B=2$ sin reemplazamiento.
\planteamiento media, varianza poblacional y enumeración de los $\binom42=6$ minibatches; comprobación con la \cref{prop:sin-reemplazo}.
\pasos
\begin{enumerate}[label=\arabic*.]
  \item Gradiente completo: $(-2+0+1+5)/4=1$.
  \item Desviaciones: $(-3,-1,0,4)$; cuadrados $(9,1,0,16)$; $\sigma^{2}=26/4=6{,}5$.
  \item Minibatches: $\{-2,0\}\to-1$; $\{-2,1\}\to-0{,}5$; $\{-2,5\}\to 1{,}5$; $\{0,1\}\to 0{,}5$; $\{0,5\}\to 2{,}5$; $\{1,5\}\to 3$.
  \item Media: $(-1-0{,}5+1{,}5+0{,}5+2{,}5+3)/6=6/6=1$.
  \item Desviaciones: $(-2;\,-1{,}5;\,0{,}5;\,-0{,}5;\,1{,}5;\,2)$; cuadrados $(4;\,2{,}25;\,0{,}25;\,0{,}25;\,2{,}25;\,4)$, que suman 13; varianza $13/6\approx 2{,}167$.
  \item Fórmula: $(6{,}5/2)\cdot(4-2)/(4-1)=3{,}25\cdot\frac23=13/6$.
\end{enumerate}
\resultado \textbf{$\nabla J=1$; $\sigma^{2}=6{,}5$; media de los minibatches $1$ (insesgado); varianza $13/6\approx 2{,}167$.}
\comprobacion el cálculo directo y la fórmula coinciden, y $2{,}167<\sigma^{2}/B=3{,}25$.
\end{solucion}

\begin{solucion}{eje:3-reemplazo}
\datos $g=(2,4,6,8)$, $N=4$, $B=2$.
\planteamiento enumerar los pares y comparar con $\sigma^{2}/B$ (\cref{teo:minibatch}) y con $\frac{\sigma^{2}}{B}\cdot\frac{N-B}{N-1}$ (\cref{prop:sin-reemplazo}).
\pasos
\begin{enumerate}[label=\arabic*.]
  \item Media $(2+4+6+8)/4=5$; $\sigma^{2}=(9+1+1+9)/4=5$.
  \item (a) Con reemplazamiento, las medias de los 16 pares ordenados toman los valores 2, 3, 4, 5, 6, 7, 8 con frecuencias 1, 2, 3, 4, 3, 2, 1.
  \item Media: $(2+6+12+20+18+14+8)/16=80/16=5$.
  \item Varianza: $(1\cdot 9+2\cdot 4+3\cdot 1+4\cdot 0+3\cdot 1+2\cdot 4+1\cdot 9)/16=40/16=2{,}5=\sigma^{2}/B$.
  \item (b) Sin reemplazamiento: $\{2,4\}\to 3$, $\{2,6\}\to 4$, $\{2,8\}\to 5$, $\{4,6\}\to 5$, $\{4,8\}\to 6$, $\{6,8\}\to 7$; media $30/6=5$.
  \item Varianza: $(4+1+0+0+1+4)/6=10/6\approx 1{,}667=2{,}5\cdot\frac23$.
\end{enumerate}
\resultado \textbf{ambos estimadores son insesgados (media 5); varianza $2{,}5$ con reemplazamiento y $5/3\approx 1{,}667$ sin él.}
\comprobacion las frecuencias del paso~2 suman $1+2+3+4+3+2+1=16$. Sin reemplazamiento no se repite ningún ejemplo y la varianza es menor.
\end{solucion}

\begin{solucion}{eje:3-suelo}
\datos $a=1$, $s^{2}=1$, $\eta\in\{0{,}1;\ 0{,}01\}$.
\planteamiento \cref{prop:suelo}: $\E[x^{2}]\to\eta s^{2}/(a(2-\eta a))$ y $\E[f]\to\eta s^{2}/(2(2-\eta a))$.
\pasos
\begin{enumerate}[label=\arabic*.]
  \item $\eta=0{,}1$: $\E[x^{2}]\to 0{,}1/1{,}9\approx 0{,}05263$; $\E[f]\to 0{,}02632$.
  \item $\eta=0{,}01$: $\E[x^{2}]\to 0{,}01/1{,}99\approx 0{,}005025$; $\E[f]\to 0{,}002513$.
  \item Cociente: $0{,}02632/0{,}002513\approx 10{,}47$.
  \item Velocidad: la media se contrae con factor $1-\eta a$, que pasa de $0{,}9$ a $0{,}99$. Para reducir $\E[x_k]$ en $e^{-1}$ hacen falta $-1/\ln(1-\eta a)$ pasos: unos $9{,}5$ con $\eta=0{,}1$ y unos $99{,}5$ con $\eta=0{,}01$.
\end{enumerate}
\resultado \textbf{el exceso de pérdida pasa de $0{,}02632$ a $0{,}002513$, unas $10{,}5$ veces menos; a cambio, la fase transitoria es unas 10 veces más lenta.} Es el compromiso que resuelve un calendario decreciente: pasos grandes al principio y pequeños al final.
\comprobacion para $\eta a\ll 1$, $\E[f]\approx\eta s^{2}/4$: $0{,}025$ y $0{,}0025$, próximos a los valores exactos.
\end{solucion}

\begin{solucion}{eje:3-sin-reemplazo}
\datos $\{I_1,\dots,I_B\}$ uniforme entre los subconjuntos de $B$ índices distintos; $X_b=\nabla\ell_{I_b}(\theta)-\nabla J(\theta)$; $d_i=\nabla\ell_i(\theta)-\nabla J(\theta)$, con $\sum_i d_i=0$ y $\frac1N\sum_i\norm{d_i}^{2}=\sigma^{2}$.
\planteamiento calcular $\E\norm{X_b}^{2}$ y $\E[X_b\T X_{b'}]$ para $b\ne b'$, y desarrollar $\Var[\gb]=\E\norm{\frac1B\sum_bX_b}^{2}$.
\pasos
\begin{enumerate}[label=\arabic*.]
  \item Por simetría, cada $I_b$ es uniforme en $\{1,\dots,N\}$, luego $\E[X_b]=\frac1N\sum_id_i=0$ (insesgadez) y $\E\norm{X_b}^{2}=\sigma^{2}$.
  \item Para $b\ne b'$, el par $(I_b,I_{b'})$ es uniforme entre los $N(N-1)$ pares ordenados de índices distintos: $\E[X_b\T X_{b'}]=\frac{1}{N(N-1)}\sum_{i\ne j}d_i\T d_j$.
  \item Como $\sum_id_i=0$: $0=\norm{\sum_id_i}^{2}=\sum_i\norm{d_i}^{2}+\sum_{i\ne j}d_i\T d_j$, luego $\sum_{i\ne j}d_i\T d_j=-N\sigma^{2}$.
  \item Por tanto $\E[X_b\T X_{b'}]=-\sigma^{2}/(N-1)$.
  \item $\Var[\gb]=\frac{1}{B^{2}}\Bigl(\sum_b\E\norm{X_b}^{2}+\sum_{b\ne b'}\E[X_b\T X_{b'}]\Bigr)=\frac{1}{B^{2}}\Bigl(B\sigma^{2}-\frac{B(B-1)\sigma^{2}}{N-1}\Bigr)$.
  \item Simplificando: $\frac{\sigma^{2}}{B}\Bigl(1-\frac{B-1}{N-1}\Bigr)=\frac{\sigma^{2}}{B}\cdot\frac{N-B}{N-1}$.
\end{enumerate}
\resultado \textbf{$\E[\gb]=\nabla J(\theta)$ y $\Var[\gb]=\frac{\sigma^{2}}{B}\cdot\frac{N-B}{N-1}$.}
\comprobacion $B=1$ da $\sigma^{2}$ y $B=N$ da~0. En el \cref{ej:minibatch}, $(3{,}5/2)\cdot(2/3)=7/6$, que es la varianza obtenida al enumerar.
\end{solucion}

\begin{solucion}{eje:3-ruido}
\datos $f(x)=x_1^{2}+2x_2^{2}$; $F(x,w)=(x_1-w_1)^{2}+2(x_2-w_2)^{2}$; $w\sim\mathcal{N}(0,s^{2}I)$.
\planteamiento derivar $F$ respecto de $x$ y separar la parte que no depende de $w$; después tomar la esperanza en $w$.
\pasos
\begin{enumerate}[label=\arabic*.]
  \item $\nabla_xF(x,w)=\bigl(2(x_1-w_1);\,4(x_2-w_2)\bigr)=\nabla f(x)+\xi$, con $\nabla f(x)=(2x_1;\,4x_2)$ y $\xi=(-2w_1;\,-4w_2)$.
  \item $\xi$ es una transformación lineal de un vector gaussiano, luego es gaussiano, de media $(0;\,0)$ y covarianza $\diag(4s^{2},16s^{2})$, porque $w_1$ y $w_2$ son independientes.
  \item $\E_w[(x_1-w_1)^{2}]=x_1^{2}-2x_1\E[w_1]+\E[w_1^{2}]=x_1^{2}+s^{2}$; del mismo modo, $\E_w[(x_2-w_2)^{2}]=x_2^{2}+s^{2}$.
  \item $\E_w[F(x,w)]=x_1^{2}+s^{2}+2(x_2^{2}+s^{2})=f(x)+3s^{2}$.
  \item La constante $3s^{2}$ no depende de $x$: $\E_w[F]$ y $f$ tienen el mismo minimizador, $x^{\ast}=0$.
\end{enumerate}
\resultado \textbf{$\nabla_xF=\nabla f+\xi$ con $\xi\sim\mathcal{N}(0,\diag(4s^{2},16s^{2}))$ y $\E_w[F]=f+3s^{2}$.} Cada paso de SGD sobre $F$ con una muestra nueva de $w$ es un paso de descenso sobre $f$ con gradiente ruidoso insesgado.
\comprobacion con $s=0$ el ruido desaparece y $F=f$.
\end{solucion}

% ---------------------------------------------------------------------
\section{Tema 4}

\begin{solucion}{eje:4-efectivo}
\datos gradiente constante $g$; $\beta\in\{0{,}9;\ 0{,}99\}$.
\planteamiento por la \cref{prop:ema}, $v_k=\frac{1-\beta^{k}}{1-\beta}g$.
\pasos
\begin{enumerate}[label=\arabic*.]
  \item Límite: $v_k\to g/(1-\beta)$. Con $\beta=0{,}9$, factor $10$; con $\beta=0{,}99$, factor $100$.
  \item El 90~\% del límite exige $1-\beta^{k}\ge 0{,}9$, es decir, $\beta^{k}\le 0{,}1$: $k\ge\ln 0{,}1/\ln 0{,}9=21{,}85$, luego $k=22$.
  \item Sobrepaso: al acercarse al mínimo, el gradiente disminuye, pero $v_k$ conserva una fracción $\beta$ de la velocidad acumulada. Con $\beta$ grande la memoria es larga y la velocidad tarda en frenarse, de modo que el iterado rebasa el mínimo.
\end{enumerate}
\resultado \textbf{factor 10 con $\beta=0{,}9$ y 100 con $\beta=0{,}99$; 22 pasos para llegar al 90~\% con $\beta=0{,}9$.}
\comprobacion $0{,}9^{22}=0{,}0985\le 0{,}1$ y $0{,}9^{21}=0{,}1094>0{,}1$.
\end{solucion}

\begin{solucion}{eje:4-pasos}
\datos $f(x)=x^{2}/2$, $f'(x)=x$, $x_0=1$, $v_0=0$, $\eta=0{,}1$, $\beta=0{,}9$.
\planteamiento aplicar la \cref{def:momentum}; el descenso de gradiente es $x\leftarrow 0{,}9x$.
\pasos
\begin{enumerate}[label=\arabic*.]
  \item $k=1$: $v_1=0+1=1$; $x_1=1-0{,}1=0{,}9$.
  \item $k=2$: $v_2=0{,}9\cdot 1+0{,}9=1{,}8$; $x_2=0{,}9-0{,}18=0{,}72$.
  \item $k=3$: $v_3=0{,}9\cdot 1{,}8+0{,}72=1{,}62+0{,}72=2{,}34$; $x_3=0{,}72-0{,}234=0{,}486$.
  \item Descenso de gradiente: $0{,}9$; $0{,}81$; $0{,}729$.
  \item Fracción recorrida: momentum $1-0{,}486=0{,}514$; descenso de gradiente $1-0{,}729=0{,}271$.
\end{enumerate}
\resultado \textbf{momentum: $0{,}9$; $0{,}72$; $0{,}486$ (51,4~\% del camino). Descenso de gradiente: $0{,}9$; $0{,}81$; $0{,}729$ (27,1~\%).}
\comprobacion en el primer paso ambos coinciden porque $v_0=0$. Por la \cref{prop:dos-pasos}, $x_3=x_2-0{,}1x_2+0{,}9(x_2-x_1)=0{,}648-0{,}162=0{,}486$.
\end{solucion}

\begin{solucion}{eje:4-optimos}
\datos $\mu=1$, $L=100$, $\kappa=100$, $\epsilon=10^{-3}$.
\planteamiento fórmulas del \cref{teo:heavy-ball} y $k=\lceil\ln\epsilon/\ln\rho\rceil$.
\pasos
\begin{enumerate}[label=\arabic*.]
  \item $\sqrt L=10$, $\sqrt\mu=1$: $\eta=4/11^{2}=4/121\approx 0{,}03306$.
  \item $\beta=\bigl(\frac{10-1}{10+1}\bigr)^{2}=\frac{81}{121}\approx 0{,}6694$; $\rho=\frac{9}{11}\approx 0{,}8182$.
  \item Iteraciones: $\ln(9/11)=-0{,}2007$; $6{,}908/0{,}2007\approx 34{,}4$, luego 35.
  \item Descenso de gradiente: $\eta^{\ast}=2/101\approx 0{,}01980$, $\rho=99/101\approx 0{,}9802$; $\ln\rho=-0{,}02000$; $6{,}908/0{,}02000\approx 345{,}4$, luego 346.
\end{enumerate}
\resultado \textbf{$\eta=4/121\approx 0{,}03306$, $\beta=81/121\approx 0{,}6694$, $\rho=9/11\approx 0{,}8182$: unas 35 iteraciones frente a 346}, una reducción en un factor próximo a $\sqrt\kappa=10$.
\comprobacion coincide con la fila $\kappa=100$ del \cref{tab:kappa}. El paso del momentum es mayor que el límite $2/L=0{,}02$ del descenso de gradiente, lo que es posible por el \cref{eje:4-estabilidad}.
\end{solucion}

\begin{solucion}{eje:4-autovalores}
\datos $\lambda=(1,\frac12,\frac14)$, $x_0=(1,1,1)$, tolerancia $10^{-2}$.
\planteamiento las coordenadas evolucionan por separado; en el descenso de gradiente, $x_k^{(i)}=(1-\eta\lambda_i)^{k}$.
\pasos
\begin{enumerate}[label=\arabic*.]
  \item (a) $\eta=1$: factores $(0;\,\frac12;\,\frac34)$. La más lenta es la tercera: $0{,}75^{k}<10^{-2}$ exige $k>\ln 0{,}01/\ln 0{,}75=16{,}01$, es decir, $k=17$.
  \item (b) $\mu=\frac14$, $L=1$: $\eta^{\ast}=2/1{,}25=1{,}6$, factores $(-0{,}6;\,0{,}2;\,0{,}6)$ y $\rho=0{,}6$. $0{,}6^{k}<10^{-2}$ exige $k>9{,}015$: $k=10$.
  \item (c) $\sqrt L=1$, $\sqrt\mu=\frac12$: $\eta=4/(1{,}5)^{2}=\frac{16}{9}\approx 1{,}778$, $\beta=\bigl(\frac{0{,}5}{1{,}5}\bigr)^{2}=\frac19$, $\rho=\frac13$.
  \item Primer paso: $v_1=\lambda\odot x_0=(1;\,\frac12;\,\frac14)$; $x_1=x_0-\frac{16}{9}v_1=\bigl(-\frac79;\,\frac19;\,\frac59\bigr)$.
  \item Segundo paso: $v_2=\frac19v_1+\lambda\odot x_1=\bigl(\frac19-\frac79;\,\frac{1}{18}+\frac{1}{18};\,\frac{1}{36}+\frac{5}{36}\bigr)=\bigl(-\frac23;\,\frac19;\,\frac16\bigr)$.
  \item $x_2=x_1-\frac{16}{9}v_2=\bigl(-\frac79+\frac{32}{27};\,\frac19-\frac{16}{81};\,\frac59-\frac{16}{54}\bigr)=\bigl(\frac{11}{27};\,-\frac{7}{81};\,\frac{7}{27}\bigr)\approx(0{,}4074;\,-0{,}08642;\,0{,}2593)$.
  \item Estimación: $(1/3)^{k}<10^{-2}$ exige $k>\ln 0{,}01/\ln(1/3)=4{,}19$, es decir, unos 5 pasos.
\end{enumerate}
\resultado \textbf{(a) 17 pasos; (b) 10 pasos; (c) $\eta=16/9$, $\beta=1/9$, $x_2=(11/27;\,-7/81;\,7/27)$ y una estimación de unos 5 pasos.} Una simulación numérica da 7 pasos: la estimación por el radio espectral es optimista porque en los autovalores extremos la raíz es doble (observación posterior al \cref{teo:heavy-ball}).
\comprobacion la forma de dos pasos (\cref{prop:dos-pasos}) en la primera coordenada da: $x_2^{(1)}=x_1^{(1)}-\eta x_1^{(1)}+\beta(x_1^{(1)}-x_0^{(1)})=-\frac79+\frac{112}{81}-\frac{16}{81}=\frac{-63+96}{81}=\frac{33}{81}=\frac{11}{27}$.
\end{solucion}

\begin{solucion}{eje:4-forma}
\datos $v_0=0$, $v_{k+1}=\beta v_k+\nabla f(\theta_k)$, $\theta_{k+1}=\theta_k-\eta v_{k+1}$.
\planteamiento eliminar $v$ entre las dos ecuaciones; para el gradiente constante, usar la \cref{prop:ema}.
\pasos
\begin{enumerate}[label=\arabic*.]
  \item De la segunda ecuación en el paso anterior, $\eta v_k=\theta_{k-1}-\theta_k$ para $k\ge 1$.
  \item $\theta_{k+1}=\theta_k-\eta\beta v_k-\eta\nabla f(\theta_k)=\theta_k-\eta\nabla f(\theta_k)+\beta(\theta_k-\theta_{k-1})$.
  \item Con $\nabla f(\theta_k)=g$: $v_{k+1}=\sum_{j=0}^{k}\beta^{k-j}g=\frac{1-\beta^{k+1}}{1-\beta}g$.
  \item $\theta_k-\theta_{k+1}=\eta v_{k+1}=\eta\frac{1-\beta^{k+1}}{1-\beta}g$.
\end{enumerate}
\resultado \textbf{$\theta_{k+1}=\theta_k-\eta\nabla f(\theta_k)+\beta(\theta_k-\theta_{k-1})$ y, con gradiente constante, $\theta_k-\theta_{k+1}=\eta\frac{1-\beta^{k+1}}{1-\beta}g$.}
\comprobacion $k=0$ da $\eta g$, el paso del descenso de gradiente; $k\to\infty$ da $\eta g/(1-\beta)$. En el \cref{ej:momentum}, $v_2=1{,}8$ y $x_1-x_2=0{,}18=0{,}1\cdot 1{,}8$.
\end{solucion}

\begin{solucion}{eje:4-estabilidad}
\datos $p(r)=r^{2}-br+\beta$ con $b=1+\beta-\eta\lambda$, $\beta\in[0,1)$, $\eta\lambda>0$.
\planteamiento distinguir raíces complejas y reales; en el caso real, usar el signo de $p(\pm 1)$.
\pasos
\begin{enumerate}[label=\arabic*.]
  \item Se calcula $p(1)=1-b+\beta=\eta\lambda$ y $p(-1)=1+b+\beta=2(1+\beta)-\eta\lambda$.
  \item \emph{Raíces complejas o doble} ($b^{2}\le 4\beta$): su módulo es $\sqrt\beta<1$. Además $\abs{b}\le 2\sqrt\beta<1+\beta$, luego $p(1)>0$ y $p(-1)>0$: la condición del enunciado se cumple.
  \item \emph{Raíces reales distintas} $r_1<r_2$: $p(r)>0$ fuera de $[r_1,r_2]$ y $p(r)<0$ en su interior. Si ambas están en $(-1,1)$, entonces $\pm 1$ están fuera y $p(\pm 1)>0$.
  \item Recíprocamente, si $p(1)>0$ y $p(-1)>0$, ningún $\pm 1$ está en $[r_1,r_2]$. Si $b\ge 2$, $p(1)=1+\beta-b\le\beta-1<0$, imposible; si $b\le-2$, $p(-1)\le\beta-1<0$, imposible. Así $\abs{b}<2$ y el vértice $b/2$ está en $(-1,1)$; como $[r_1,r_2]$ contiene el vértice y no contiene $\pm 1$, está contenido en $(-1,1)$.
  \item En todos los casos: las raíces tienen módulo menor que~1 si y solo si $p(1)>0$ y $p(-1)>0$, es decir, $0<\eta\lambda<2(1+\beta)$.
\end{enumerate}
\resultado \textbf{el método es estable si y solo si $0<\eta\lambda<2(1+\beta)$}, frente a $0<\eta\lambda<2$ en el descenso de gradiente: con $\beta=0{,}9$ el paso puede ser hasta $1{,}9$ veces mayor.
\comprobacion con $\beta=0$ se recupera $0<\eta\lambda<2$. En el \cref{eje:4-optimos}, $\eta L=400/121\approx 3{,}306<2(1+0{,}6694)=3{,}339$.
\end{solucion}

% ---------------------------------------------------------------------
\section{Tema 5}

\begin{solucion}{eje:5-escala-test}
\datos Adam con $m_0=v_0=0$; la pérdida se multiplica por $c=1000$, de modo que $g$ pasa a ser $cg$.
\planteamiento calcular el primer paso de Adam en función de $g$ (\cref{prop:escala}).
\pasos
\begin{enumerate}[label=\arabic*.]
  \item $m_1=(1-\beta_1)g$ y $v_1=(1-\beta_2)g^{2}$.
  \item $\hat m_1=g$ y $\hat v_1=g^{2}$, luego $\sqrt{\hat v_1}=\abs{g}$.
  \item Paso: $\eta g/(\abs{g}+\varepsilon)\approx\eta\,\sign(g)$.
  \item Con $cg$: $\eta cg/(c\abs{g}+\varepsilon)\approx\eta\,\sign(g)$; el factor $c$ se cancela.
  \item En el SGD, el paso $\eta g$ pasa a ser $\eta cg$: se multiplica por 1000.
\end{enumerate}
\resultado \textbf{falsa}: el primer paso de Adam mide $\eta$ en ambos casos, salvo el efecto despreciable de $\varepsilon$; en el SGD, en cambio, sí se multiplicaría por 1000.
\comprobacion con $\eta=0{,}01$, $\varepsilon=10^{-8}$ y $g=0{,}5$, el paso es $0{,}01\cdot 0{,}5/(0{,}5+10^{-8})\approx 0{,}0099999998$; con $g=500$, $\approx 0{,}01$ con diez cifras.
\end{solucion}

\begin{solucion}{eje:5-rmsprop}
\datos RMSProp con $s_0=0$, $\varepsilon\approx 0$, gradiente inicial $g_1\ne 0$; $\rho\in\{0{,}9;\ 0{,}999\}$.
\planteamiento $s_1=(1-\rho)g_1^{2}$ y paso $\eta g_1/\sqrt{s_1}$.
\pasos
\begin{enumerate}[label=\arabic*.]
  \item Paso: $\eta g_1/(\sqrt{1-\rho}\,\abs{g_1})=\eta\,\sign(g_1)/\sqrt{1-\rho}$.
  \item $\rho=0{,}9$: $1/\sqrt{0{,}1}\approx 3{,}162$.
  \item $\rho=0{,}999$: $1/\sqrt{0{,}001}\approx 31{,}62$.
\end{enumerate}
\resultado \textbf{el primer paso mide unas $3{,}16\,\eta$ con $\rho=0{,}9$ y $31{,}6\,\eta$ con $\rho=0{,}999$.} El estimador $s_t$ está sesgado hacia~0 por la inicialización, lo que agranda los primeros pasos justo cuando el gradiente es menos fiable; la corrección de sesgo de Adam divide $v_t$ por $1-\beta_2^{t}$ y elimina el efecto.
\comprobacion con corrección, $\hat s_1=s_1/(1-\rho)=g_1^{2}$ y el paso vuelve a ser $\eta$.
\end{solucion}

\begin{solucion}{eje:5-adam}
\datos $\eta=0{,}01$, $\theta_0=1$, $\beta_1=0{,}9$, $\beta_2=0{,}999$, $\varepsilon\approx 0$, $g_1=0{,}5$, $g_2=-0{,}1$.
\planteamiento aplicar la \cref{def:adam} paso a paso.
\pasos
\begin{enumerate}[label=\arabic*.]
  \item $t=1$: $m_1=0{,}1\cdot 0{,}5=0{,}05$; $v_1=0{,}001\cdot 0{,}25=0{,}00025$.
  \item $\hat m_1=0{,}05/0{,}1=0{,}5$; $\hat v_1=0{,}00025/0{,}001=0{,}25$; $\sqrt{\hat v_1}=0{,}5$.
  \item Paso $0{,}01\cdot 0{,}5/0{,}5=0{,}01$; $\theta_1=0{,}99$.
  \item $t=2$: $m_2=0{,}9\cdot 0{,}05+0{,}1\cdot(-0{,}1)=0{,}045-0{,}01=0{,}035$.
  \item $v_2=0{,}999\cdot 0{,}00025+0{,}001\cdot 0{,}01=0{,}00024975+0{,}00001=0{,}00025975$.
  \item Correcciones: $1-0{,}9^{2}=0{,}19$ y $1-0{,}999^{2}=0{,}001999$.
  \item $\hat m_2=0{,}035/0{,}19=0{,}1842$; $\hat v_2=0{,}00025975/0{,}001999=0{,}1299$; $\sqrt{\hat v_2}=0{,}3605$.
  \item Paso $0{,}01\cdot 0{,}1842/0{,}3605=0{,}005110$; $\theta_2=0{,}99-0{,}005110=0{,}98489$.
\end{enumerate}
\resultado \textbf{$\theta_1=0{,}99$ y $\theta_2\approx 0{,}9849$.}
\comprobacion aunque $g_2<0$, $\theta$ sigue bajando porque $m_2=0{,}035>0$ conserva la inercia del primer gradiente; el paso de $t=2$ es menor que $\eta$, como corresponde a un cambio de signo del gradiente.
\end{solucion}

\begin{solucion}{eje:5-precond}
\datos $f(x)=0{,}1x_1^{2}+2x_2^{2}$; $R$ rotación de ángulo $\pi/4$, $R=\frac{1}{\sqrt2}\begin{pmatrix}1&-1\\1&1\end{pmatrix}$; $g(y)=f(R\T y)$.
\planteamiento \cref{prop:precond}: el hessiano preacondicionado es $D^{-1/2}QD^{-1/2}$, semejante a $D^{-1}Q$; para $g$, el hessiano es $RQR\T$ por la regla de la cadena.
\pasos
\begin{enumerate}[label=\arabic*.]
  \item $Q=\nabla^{2}f=\diag(0{,}2;\,4)$, $\kappa=4/0{,}2=20$.
  \item $D=\diag(0{,}2;\,4)=Q$: $D^{-1/2}QD^{-1/2}=I$, $\kappa=1$.
  \item $\nabla^{2}g=RQR\T$. Con $a=0{,}2$ y $b=4$: $RQR\T=\frac12\begin{pmatrix}a+b&a-b\\a-b&a+b\end{pmatrix}=\begin{pmatrix}2{,}1&-1{,}9\\-1{,}9&2{,}1\end{pmatrix}$.
  \item Diagonal: $D'=\diag(2{,}1;\,2{,}1)=2{,}1\,I$. Preacondicionar por $D'$ divide el hessiano por $2{,}1$: autovalores $0{,}2/2{,}1\approx 0{,}09524$ y $4/2{,}1\approx 1{,}905$, con $\kappa=20$.
  \item (c) El preacondicionador diagonal elimina el mal condicionamiento cuando los ejes principales coinciden con los coordenados, pero no corrige nada si están girados.
\end{enumerate}
\resultado \textbf{(a) $\kappa=20$, y $\kappa=1$ tras preacondicionar; (b) $\nabla^{2}g=\begin{pmatrix}2{,}1&-1{,}9\\-1{,}9&2{,}1\end{pmatrix}$, diagonal $2{,}1\,I$ y $\kappa=20$ tras preacondicionar.} Los métodos adaptativos diagonales dependen de la base: funcionan bien cuando las escalas distintas están alineadas con los parámetros.
\comprobacion la rotación conserva los autovalores: traza $4{,}2=0{,}2+4$ y determinante $2{,}1^{2}-1{,}9^{2}=0{,}8=0{,}2\cdot 4$.
\end{solucion}

\begin{solucion}{eje:5-memoria}
\datos $N=7\cdot 10^{9}$ parámetros; 4 bytes por número; se guardan parámetros, gradientes, $m$ y $v$.
\planteamiento cada vector ocupa $4N$ bytes.
\pasos
\begin{enumerate}[label=\arabic*.]
  \item Un vector: $4\cdot 7\cdot 10^{9}=2{,}8\cdot 10^{10}$ bytes $=28$~GB.
  \item Cuatro vectores: $4\cdot 28=112$~GB.
  \item Estado del optimizador ($m$ y $v$): $2\cdot 28=56$~GB.
  \item Fracción: $56/112=0{,}5$.
\end{enumerate}
\resultado \textbf{112~GB en total, de los que 56~GB (el 50~\%) son el estado de Adam}, sin contar las activaciones.
\comprobacion el SGD sin momentum no tiene estado y necesitaría 56~GB; con momentum, 84~GB. Las técnicas de reparto del estado entre dispositivos se estudian en el módulo~17.
\end{solucion}

\begin{solucion}{eje:5-invariancia}
\datos Adam con $\varepsilon=0$; gradientes $g_t$ sustituidos por $g'_t=cg_t$.
\planteamiento por inducción sobre $t$: si los iterados coinciden hasta $t-1$, los gradientes de ambas ejecuciones en $\theta_{t-1}$ están en razón $c$, y se comprueba cómo escalan $m_t$, $v_t$ y el cociente.
\pasos
\begin{enumerate}[label=\arabic*.]
  \item Hipótesis: $\theta'_s=\theta_s$, $m'_s=cm_s$ y $v'_s=c^{2}v_s$ para $s\le t-1$ (cierto para $s=0$, con $m_0=v_0=0$).
  \item Entonces $g'_t=cg_t$, $m'_t=\beta_1cm_{t-1}+(1-\beta_1)cg_t=cm_t$ y $v'_t=\beta_2c^{2}v_{t-1}+(1-\beta_2)c^{2}g_t^{2}=c^{2}v_t$.
  \item $\hat m'_t/\sqrt{\hat v'_t}=c\hat m_t/(\abs{c}\sqrt{\hat v_t})=\hat m_t/\sqrt{\hat v_t}$ si $c>0$. Luego $\theta'_t=\theta_t$.
  \item Si $c<0$, el cociente cambia de signo: los iterados se mueven en sentido contrario (se asciende por $f$). No hay invariancia.
  \item SGD con momentum: $v'_k=cv_k$ y el paso es $\eta cv_k$: se multiplica por $c$. No hay invariancia.
\end{enumerate}
\resultado \textbf{los iterados de Adam son invariantes frente a $g\mapsto cg$ con $c>0$}, pero no con $c<0$, y el SGD con momentum no tiene esta propiedad.
\comprobacion con $c>0$ y $\varepsilon>0$, el cociente es $c\hat m_t/(c\sqrt{\hat v_t}+\varepsilon)$, que solo coincide si $\varepsilon\ll c\sqrt{\hat v_t}$; de ahí que $\varepsilon$ se elija muy pequeño.
\end{solucion}

\begin{solucion}{eje:5-sin-correccion}
\datos $\hat m_t=m_t/(1-\beta_1^{t})$, $\hat v_t=v_t/(1-\beta_2^{t})$, $\varepsilon\to 0$, $\beta_1=0{,}9$, $\beta_2=0{,}999$.
\planteamiento sacar los factores de corrección fuera del cociente.
\pasos
\begin{enumerate}[label=\arabic*.]
  \item $\eta\dfrac{\hat m_t}{\sqrt{\hat v_t}}=\eta\,\dfrac{m_t/(1-\beta_1^{t})}{\sqrt{v_t}/\sqrt{1-\beta_2^{t}}}=\eta\,\dfrac{\sqrt{1-\beta_2^{t}}}{1-\beta_1^{t}}\cdot\dfrac{m_t}{\sqrt{v_t}}=\eta_t\,\dfrac{m_t}{\sqrt{v_t}}$.
  \item $t=1$: $\sqrt{0{,}001}/0{,}1=0{,}03162/0{,}1=0{,}3162$.
  \item $t=2$: $\sqrt{0{,}001999}/0{,}19=0{,}04471/0{,}19=0{,}2353$.
  \item $t=10$: $1-0{,}999^{10}=0{,}009955$, de raíz $0{,}09978$; $1-0{,}9^{10}=0{,}6513$; cociente $0{,}1532$.
  \item $t=1000$: $1-0{,}999^{1000}=0{,}6323$, de raíz $0{,}7952$; $1-0{,}9^{1000}\approx 1$; cociente $0{,}7952$.
  \item Paso de $t=2$: $m_2=0{,}035$ y $\sqrt{v_2}=\sqrt{0{,}00025975}=0{,}01612$; $0{,}01\cdot 0{,}2353\cdot 0{,}035/0{,}01612=0{,}005109$.
\end{enumerate}
\resultado \textbf{$\eta_t/\eta=0{,}3162$; $0{,}2353$; $0{,}1532$; $0{,}7952$ para $t=1$, 2, 10 y 1000.} La corrección equivale a una tasa que crece hacia $\eta$ durante los primeros miles de pasos: un calentamiento implícito.
\comprobacion el paso de $t=2$ coincide con el del \cref{eje:5-adam} ($0{,}00511$) salvo redondeo. Cuando $t\to\infty$, ambos factores tienden a~1 y $\eta_t\to\eta$.
\end{solucion}

% ---------------------------------------------------------------------
\section{Tema 6}

\begin{solucion}{eje:6-warmup}
\datos calendario de calentamiento lineal y decaimiento coseno (\cref{def:coseno}).
\planteamiento identificar qué falla al principio y al final del entrenamiento con $\eta$ constante.
\pasos
\begin{enumerate}[label=\arabic*.]
  \item \emph{Al principio.} Los gradientes son grandes e irregulares y los momentos de Adam se estiman con muy pocos gradientes. Una tasa grande desde el primer paso puede hacer divergir la pérdida (\cref{teo:gd-cuadratica}: el paso debe ser menor que $2/\lambda_{\max}$). El calentamiento sube $\eta$ gradualmente; el \cref{eje:5-sin-correccion} muestra que la corrección de sesgo de Adam ya produce un efecto parecido.
  \item \emph{Al final.} Con gradientes estocásticos, el exceso de pérdida estacionario es proporcional a $\eta$ (\cref{prop:suelo}). El decaimiento coseno reduce $\eta$ suavemente y rebaja ese suelo de ruido.
\end{enumerate}
\resultado \textbf{el calentamiento aporta estabilidad inicial y el decaimiento, precisión final.}
\comprobacion en el \cref{eje:3-suelo}, dividir $\eta$ por 10 divide el exceso de pérdida por $10{,}5$.
\end{solucion}

\begin{solucion}{eje:6-calendario}
\datos $\eta_{\max}=6\cdot 10^{-4}$, $\eta_{\min}=6\cdot 10^{-5}$, $W=1000$, $T=11\,000$.
\planteamiento \cref{def:coseno}; progreso del tramo coseno $(t-1000)/10\,000$; amplitud $\frac12(\eta_{\max}-\eta_{\min})=2{,}7\cdot 10^{-4}$.
\pasos
\begin{enumerate}[label=\arabic*.]
  \item $t=500<W$: $\eta=6\cdot 10^{-4}\cdot 500/1000=3\cdot 10^{-4}$.
  \item $t=6000$: progreso $0{,}5$; $\cos(\pi/2)=0$; $\eta=6\cdot 10^{-5}+2{,}7\cdot 10^{-4}=3{,}3\cdot 10^{-4}$.
  \item $t=8500$: progreso $0{,}75$; $\cos(0{,}75\pi)=-0{,}7071$; $1-0{,}7071=0{,}2929$.
  \item $\eta=6\cdot 10^{-5}+2{,}7\cdot 10^{-4}\cdot 0{,}2929=6\cdot 10^{-5}+7{,}908\cdot 10^{-5}=1{,}391\cdot 10^{-4}$.
\end{enumerate}
\resultado \textbf{$\eta(500)=3\cdot 10^{-4}$; $\eta(6000)=3{,}3\cdot 10^{-4}$; $\eta(8500)\approx 1{,}391\cdot 10^{-4}$.}
\comprobacion en el punto medio, $(\eta_{\max}+\eta_{\min})/2=(6+0{,}6)\cdot 10^{-4}/2=3{,}3\cdot 10^{-4}$ (\cref{prop:coseno}).
\end{solucion}

\begin{solucion}{eje:6-adamw}
\datos $\eta=10^{-3}$, $\lambda=0{,}1$, $\theta=1$, $\hat v=10^{-4}$, $\varepsilon\approx 0$.
\planteamiento con $L^{2}$, el término $\lambda\theta$ queda dividido por $\sqrt{\hat v}$ (\cref{obs:adam-l2}); en AdamW se resta $\eta\lambda\theta$ (\cref{def:adamw}).
\pasos
\begin{enumerate}[label=\arabic*.]
  \item $\sqrt{\hat v}=10^{-2}$.
  \item Adam con $L^{2}$: $\eta\lambda\theta/\sqrt{\hat v}=10^{-3}\cdot 0{,}1/10^{-2}=10^{-2}$.
  \item AdamW: $\eta\lambda\theta=10^{-3}\cdot 0{,}1=10^{-4}$.
  \item Cociente: $10^{-2}/10^{-4}=100$.
\end{enumerate}
\resultado \textbf{Adam con $L^{2}$ encoge $0{,}01$ ($\theta\to 0{,}99$) y AdamW $0{,}0001$ ($\theta\to 0{,}9999$); el cociente es 100.} El parámetro de gradientes pequeños se regulariza 100 veces más con $L^{2}$.
\comprobacion con $\hat v=1$ ambos métodos coinciden en $10^{-4}$, como en el \cref{tab:adamw}.
\end{solucion}

\begin{solucion}{eje:6-semivida}
\datos (a) factor por paso $1-10^{-4}$; (b) $\eta=3\cdot 10^{-4}$, $\lambda=0{,}1$, $10^{5}$ pasos.
\planteamiento tras $k$ pasos de decaimiento puro, el peso queda multiplicado por $(1-\eta\lambda)^{k}$.
\pasos
\begin{enumerate}[label=\arabic*.]
  \item (a) $(1-10^{-4})^{k}=\frac12$, luego $k=\ln 0{,}5/\ln(1-10^{-4})=-0{,}6931/(-1{,}0001\cdot 10^{-4})\approx 6931$.
  \item (b) $\eta\lambda=3\cdot 10^{-5}$; $\ln\bigl((1-3\cdot 10^{-5})^{10^{5}}\bigr)=10^{5}\ln(1-3\cdot 10^{-5})\approx-3{,}0000$.
  \item $e^{-3}\approx 0{,}04979$.
\end{enumerate}
\resultado \textbf{(a) unos 6931 pasos; (b) el peso queda multiplicado por $\approx 0{,}0498$.}
\comprobacion para $x$ pequeño, $(1-x)^{k}\approx e^{-kx}$: en (a), $e^{-6931\cdot 10^{-4}}=e^{-0{,}6931}=0{,}5$.
\end{solucion}

\begin{solucion}{eje:6-ridge}
\datos $J_\lambda(w)=\frac{a}{2}(w-c)^{2}+\frac{\lambda}{2}w^{2}$, $a>0$, $\lambda\ge 0$; caso numérico $a=2$, $c=3$, $\lambda=2$.
\planteamiento anular la derivada; $J_\lambda''=a+\lambda>0$ garantiza que es el mínimo (\cref{teo:opt-convexa}).
\pasos
\begin{enumerate}[label=\arabic*.]
  \item $J_\lambda'(w)=a(w-c)+\lambda w$.
  \item $J_\lambda'(w^{\ast})=0\Rightarrow(a+\lambda)w^{\ast}=ac\Rightarrow w^{\ast}=ac/(a+\lambda)$.
  \item Paso de descenso: $w-\eta J_\lambda'(w)=w-\eta\lambda w-\eta a(w-c)=(1-\eta\lambda)w-\eta a(w-c)$.
  \item Números: $w^{\ast}=2\cdot 3/(2+2)=1{,}5$.
  \item Convergencia: $J_\lambda$ es una cuadrática de curvatura $a+\lambda$; por el \cref{ej:gd-1d}, converge si y solo si $0<\eta<2/(a+\lambda)$, que con los datos es $0<\eta<0{,}5$.
\end{enumerate}
\resultado \textbf{$w^{\ast}=ac/(a+\lambda)=1{,}5$; el paso de descenso es $(1-\eta\lambda)w-\eta a(w-c)$ y converge para $0<\eta<2/(a+\lambda)$.}
\comprobacion $J_\lambda'(1{,}5)=2\cdot(1{,}5-3)+2\cdot 1{,}5=-3+3=0$. Con $\lambda=0$ se recupera $w^{\ast}=c=3$, y el factor $a/(a+\lambda)=1/2$ es el de la \cref{prop:ridge}.
\end{solucion}

\begin{solucion}{eje:6-polinomico}
\datos $\eta_t=\eta_0t^{-\alpha}$, $t\ge 1$, $\eta_0>0$, $\alpha>0$.
\planteamiento criterio de las series $\sum t^{-p}$ (convergen si y solo si $p>1$) y comparación con integrales para acotar sumas de funciones decrecientes.
\pasos
\begin{enumerate}[label=\arabic*.]
  \item (a) $\sum_t\eta_t=\eta_0\sum_tt^{-\alpha}$ diverge si y solo si $\alpha\le 1$.
  \item $\sum_t\eta_t^{2}=\eta_0^{2}\sum_tt^{-2\alpha}$ converge si y solo si $2\alpha>1$, es decir, $\alpha>\frac12$.
  \item Ambas condiciones a la vez: $\frac12<\alpha\le 1$.
  \item (b) $\alpha=\frac12$. Como $x\mapsto x^{-1/2}$ es decreciente, $t^{-1/2}\ge\int_t^{t+1}x^{-1/2}\,dx$; sumando, $\sum_{t=1}^{T}t^{-1/2}\ge\int_1^{T+1}x^{-1/2}\,dx=2(\sqrt{T+1}-1)$.
  \item $\eta_t^{2}=\eta_0^{2}/t$ y, para $t\ge 2$, $1/t\le\int_{t-1}^{t}dx/x$; así $\sum_{t=1}^{T}1/t\le 1+\int_1^{T}dx/x=1+\ln T$.
  \item Cota del \cref{teo:sgd-convexo}: $\dfrac{r^{2}+G^{2}\eta_0^{2}(1+\ln T)}{4\eta_0(\sqrt{T+1}-1)}$. El numerador crece como $\ln T$ y el denominador como $\sqrt T$.
\end{enumerate}
\resultado \textbf{(a) las condiciones de Robbins--Monro se cumplen si y solo si $\frac12<\alpha\le 1$; (b) con $\alpha=\frac12$ la cota es $O(\ln T/\sqrt T)$}, casi la velocidad $O(1/\sqrt T)$ del paso constante óptimo (\cref{cor:sgd-constante}), sin necesidad de conocer $T$ de antemano.
\comprobacion con $T=100$: $\sum_{t=1}^{100}t^{-1/2}\approx 18{,}59\ge 2(\sqrt{101}-1)\approx 18{,}10$, y $\sum_{t=1}^{100}1/t\approx 5{,}187\le 1+\ln 100\approx 5{,}605$.
\end{solucion}
